%% kapa: exam=1 round=23 year=2012 period=2 form=A subject=appraisal_law subject_ko=부동산관계법규 seq=13 number=53
%% source: Q-Net 감정평가사 기출문제 (cst00309) artlSeq=5052580 file=raw/1008/5052580/2012년 제23회 감정평가사 제1차시험 기출문제(2교시 A형).pdf
%% conversion: kordoc 4.24.51 -> tools/gen1.py (auto); review=proofread-sonnet; flags=none
\begin{question}{53}
국토의 계획 및 이용에 관한 법령상 행정청인 도시계획시설사업의 시행자가 사업시행과 관련하여 타인의 토지를 출입하는 경우 등에 관한 설명으로 옳지 않은 것은?
\begin{choices}
  \item 타인의 토지에 출입하려면 출입하려는 날의 3일 전까지 그 토지의 소유자ㆍ점유자 또는 관리인에게 그 일시와 장소를 알려야 한다.
  \item 타인의 토지를 임시통로로 일시사용하려면 그 토지의 소유자ㆍ점유자 또는 관리인의 동의를 받아야 한다.
  \item 타인의 토지를 재료적치장으로 일시사용하고자 할 때, 그 토지의 소유자ㆍ점유자 또는 관리인의 주소가 불분명하여 동의를 받을 수 없는 경우에는 미리 관할 시장 또는 군수의 허가를 받아야 한다.
  \item 장애물을 제거하려면 제거하려는 날의 3일 전까지 장애물의 소유자ㆍ점유자 또는 관리인에게 알려야 한다.
  \item 일출 전이나 일몰 후에는 그 토지 점유자의 승낙 없이 택지나 담장 또는 울타리로 둘러싸인 타인의 토지에 출입할 수 없다.
\end{choices}
\end{question}
\answer{3}
